\section*{Chapitre \ref{ch:b2:ellingham} --- Diagrammes d'Ellingham et métallurgie}

\begin{solution}{exo:b2:ellingham:1}
\ce{4Cu + O2 -> 2Cu2O}~; \ce{4/3Al + O2 -> 2/3Al2O3}~; \ce{2C + O2 -> 2CO}~;
\ce{C + O2 -> CO2}.
\end{solution}

\begin{solution}{exo:b2:ellingham:2}
$\Delta_r H^\circ = 2(-350.46) = \qty{-700.92}{kJ}$~; $\Delta_r S^\circ = 2(43.65)
- 2(41.63) - 205.15 = \qty{-201.1}{J/K}$. Donc $\Delta_r G^\circ = -700.9 +
0.2011\,T$ (\unit{kJ}, $T$ en \unit{K}), valable jusqu'à la température de fusion du zinc.
\end{solution}

\begin{solution}{exo:b2:ellingham:3}
À \qty{1500}{K}, la droite de \ce{Cr2O3} est vers \qty{-495}{kJ}. Au-dessous, donc
capables de réduire l'oxyde de chrome~: silicium, titane, aluminium, magnésium,
calcium. Au-dessus, incapables~: cuivre, fer, zinc. Le carbone, à \qty{-488}{kJ}, se trouve
juste au-dessus~: il ne réduit l'oxyde de chrome qu'à une température un peu plus élevée,
et donne alors un chrome chargé de carbure --- l'une des raisons pour lesquelles le chrome destiné
aux alliages exempts de carbone est produit par \omterm{def:b2:ellingham:aluminothermy}{aluminothermie}.
\end{solution}

\begin{solution}{exo:b2:ellingham:4}
La pente est $-\Delta_r S^\circ$, gouvernée par la variation de la quantité de gaz.
\ce{C + O2 -> CO2}~: une mole de gaz en donne une, $\Delta_r S^\circ \approx 0$,
droite horizontale. \ce{2C + O2 -> 2CO}~: une en donne deux, $\Delta_r S^\circ > 0$, la
droite descend. \ce{2CO + O2 -> 2CO2}~: trois en donnent deux, $\Delta_r S^\circ < 0$, la
droite monte, comme celles des métaux.
\end{solution}

\begin{solution}{exo:b2:ellingham:5}
$\Delta_r G^\circ(\qty{600}{K}) = -181.6 + 600 \times 0.2165 = \qty{-51.7}{kJ}$,
$p_{\ce{O2}} = p^\circ\exp(-51\,710/(8.314 \times 600)) = \qty{3.2e-5}{bar}$. L'air
en contient \qty{0.21}{bar}, bien davantage~: le mercure est oxydé à \qty{600}{K} (lentement
--- c'est ainsi que l'oxyde fut préparé pour la première fois), et l'oxyde ne se décompose qu'au-dessus
d'environ \qty{730}{K}.
\end{solution}

\begin{solution}{exo:b2:ellingham:6}
$\Delta_r G^\circ = -1582.3 - (-743.5) = \qty{-838.8}{kJ}$. Les réactifs sont
deux solides en contact seulement à la surface de leurs grains, et la réaction
a une énergie d'activation élevée~: il faut l'amorcer localement (un ruban de magnésium)~;
ensuite, la chaleur qu'elle libère l'entretient.
\end{solution}

\begin{solution}{exo:b2:ellingham:7}
La droite de \ce{2C + O2 -> 2CO} coupe celle de \ce{2Fe + O2 -> 2FeO} vers
\qty{1040}{K}~; au-dessus, \ce{FeO + C -> Fe + CO} a une \omterm{def:b2:reaction-free-energy:reaction-gibbs}{enthalpie libre standard de réaction}
négative.
\end{solution}

\begin{solution}{exo:b2:ellingham:8}
$\Delta_r G^\circ = 2(-209.1) - (-396.0) = \qty{-22.2}{kJ}$, $K^\circ =
\exp(22\,200/(8.314 \times 1100)) = 11.3$. Alors $x^2/(1 - x) = 11.3$, $x =
0.92$. Par refroidissement à \qty{700}{K}, l'équilibre recule vers \ce{CO2}
($x = 0.016$)~: le monoxyde de carbone se dismute, \ce{2CO -> C + CO2}, et
de la suie se dépose --- lentement, sauf si une surface comme le fer catalyse la réaction.
\end{solution}

\begin{solution}{exo:b2:ellingham:9}
La droite de l'eau (trois moles de gaz en donnent deux) monte moins fortement que la
droite \ce{CO}/\ce{CO2}, elle aussi trois pour deux mais avec une plus grande perte d'entropie. Les
deux se coupent vers \qty{1100}{K}~: au-dessous, la droite \ce{CO}/\ce{CO2} est la plus basse et
le monoxyde de carbone est le meilleur réducteur~; au-dessus, la droite de l'eau est la plus basse et
le dihydrogène l'emporte. C'est le même énoncé que le \omterm{def:b2:equilibrium-shifts:shift}{déplacement de l'équilibre}
du gaz à l'eau \ce{CO + H2O <=> CO2 + H2} vers la gauche à haute température.
\end{solution}

\begin{solution}{exo:b2:ellingham:10}
Par mole de \ce{O2}, c'est-à-dire pour \ce{2MgO + Si -> SiO2 + 2Mg(g)}~: $\Delta_r
G^\circ = -643.7 - (-845.5) = \qty{+201.8}{kJ}$. Le silicium et la silice étant
d'activité 1, $\Delta_r G = \Delta_r G^\circ + RT\ln(p_{\ce{Mg}}/p^\circ)^2$, négative
lorsque $p_{\ce{Mg}} < p^\circ\exp(-201\,800/(2 \times 8.314 \times 1500)) =
\qty{3e-4}{bar}$. Un four sous vide pompe la vapeur de magnésium à mesure
qu'elle se forme et la condense sur une surface froide, de sorte que sa pression n'atteint jamais
cette valeur~; de la chaux ajoutée à la charge fixe la silice et abaisse son activité,
ce qui aide encore.
\end{solution}

\begin{solution}{exo:b2:ellingham:11}
Quatre espèces, une réaction, trois phases (deux solides et le gaz)~: $v = 4 - 1 +
2 - 3 = 2$. À température et pression données, le rapport $p_{\ce{CO}}/p_{\ce{CO2}}$
est fixé~: le gaz ne peut pas céder tout son monoxyde de carbone à l'oxyde de fer,
et le gaz de gueulard en contient encore beaucoup --- on le brûle pour préchauffer
l'air.
\end{solution}

\begin{solution}{exo:b2:ellingham:12}
$\log K = 2(0.34 + 0.41)/0.059 = 25.4$, $K = \num{3e25}$. Avec
$[\ce{Fe^2+}] = \qty{1.0}{mol/L}$, $[\ce{Cu^2+}] = 1/K \approx
\qty{4e-26}{mol/L}$~: le cuivre est entièrement éliminé.
\end{solution}

\begin{solution}{pb:b2:ellingham:1}
\textbf{1.} $\Delta_r H^\circ = \qty{-700.92}{kJ}$, $\Delta_r S^\circ =
\qty{-201.1}{J/K}$.
\textbf{2.} $\Delta_r G^\circ = -700.9 + 0.2011\,T$ (\unit{kJ}), au-dessous de
\qty{692.7}{K}.
\textbf{3.} $\Delta_{\text{fus}}S = 7320/692.7 = \qty{10.6}{J/(K.mol)}$~;
$\Delta_{\text{vap}}S = 115\,300/1180.2 = \qty{97.7}{J/(K.mol)}$.
\textbf{4.} Zinc liquide~: $\Delta_r H^\circ = -700.92 - 2(7.32) = \qty{-715.6}{kJ}$,
$\Delta_r S^\circ = -201.1 - 2(10.6) = \qty{-222.2}{J/K}$~: $\Delta_r G^\circ =
-715.6 + 0.2222\,T$.
\textbf{5.} Zinc gazeux~: $\Delta_r H^\circ = -715.6 - 2(115.3) = \qty{-946.2}{kJ}$,
$\Delta_r S^\circ = -222.2 - 2(97.7) = \qty{-417.7}{J/K}$~: $\Delta_r G^\circ =
-946.2 + 0.4177\,T$.
\textbf{6.} Pentes $0.201$, $0.222$ et $0.418$ \unit{kJ/K}. Au-dessus de la température
d'ébullition, la réaction consomme trois moles de gaz au lieu d'une~: la perte
d'entropie, donc la pente, double.
\textbf{7.} À \qty{692.7}{K}~: $-700.9 + 139.3 = -715.6 + 153.9 = \qty{-561.6}{kJ}$~;
à \qty{1180.2}{K}~: $-715.6 + 262.3 = -946.2 + 492.9 = \qty{-453.3}{kJ}$. Les
droites se raccordent, comme il se doit.
\textbf{8.} Pente $(-453.0 + 418.2)/200 = \qty{-0.174}{kJ/K}$, $a = -418.2 +
0.174 \times 1100 = \qty{-226.8}{kJ}$~: $\Delta_r G^\circ = -226.8 - 0.174\,T$.
\textbf{9.} $\Delta_r S^\circ = \qty{+174}{J/K}$~: une mole de gaz en donne deux.
\textbf{10.} \ce{2ZnO + 2C -> 2Zn(g) + 2CO}, $\Delta_r G^\circ = (-226.8 - 0.174\,T)
- (-946.2 + 0.4177\,T) = 719.4 - 0.592\,T$ (\unit{kJ}).
\textbf{11.} Nulle pour $T = 719.4/0.592 = \qty{1215}{K}$.
\textbf{12.} Au-dessus de sa température d'ébullition, le zinc se forme à l'état de vapeur~: la réaction produit
du gaz à partir de solides, son entropie est grande et positive, et la vapeur quitte la
charge, ce qui sépare le métal du minerai et du coke par distillation.
\textbf{13.} Quatre espèces, une réaction, une condition particulière ($p_{\ce{Zn}} =
p_{\ce{CO}}$), trois phases~: $v = 4 - 1 - 1 + 2 - 3 = 1$. Fixer la pression
fixe la température d'équilibre.
\textbf{14.} Une mole de vapeur de zinc par mole de monoxyde de carbone~:
$p_{\ce{Zn}} = p_{\ce{CO}} = \qty{0.5}{bar}$.
\textbf{15.} Par mole de \ce{ZnO}, $\Delta_r G^\circ = \frac12(719.4 - 0.592 \times
1300) = \qty{-25.1}{kJ/mol}$, $K^\circ = \exp(25\,100/(8.314 \times 1300)) = 10$~;
$Q = 0.25 < K^\circ$~: la réduction se produit.
\textbf{16.} Juste au-dessus du croisement, la réduction est déjà favorisée (avec
\qty{0.5}{bar} de chaque gaz, dès environ \qty{1170}{K})~; chauffer davantage coûte du combustible,
use les cornues d'argile et ne produit pas plus de zinc, puisque la réaction est
totale de toute façon.
\textbf{17.} \ce{Zn(g) + CO2 -> ZnO + CO}.
\textbf{18.} La droite de l'oxyde de zinc est au-dessous de la droite \ce{CO}/\ce{CO2}
($-445$ contre \qty{-357}{kJ} à \qty{1200}{K})~: la vapeur de zinc réduit le dioxyde
de carbone et redevient oxyde. Dans la cornue, le carbone chaud détruit tout
\ce{CO2} (Boudouard)~; dans le condenseur plus froid, rien ne le fait, et le \ce{CO2} --- issu
d'entrées d'air, ou de \ce{2CO -> C + CO2} au refroidissement --- transforme le zinc en
une poudre grise oxydée.
\textbf{19.} Un refroidissement lent laisse longtemps la vapeur dans le domaine où
elle se réoxyde, et donne une fine poussière au lieu de métal liquide~; une condensation
rapide en liquide limite les deux effets.
\textbf{20.} $10^6/65.38 = \qty{1.53e4}{mol}$ de zinc, autant de carbone~:
$1.53 \times 10^4 \times 12.011 = \qty{184}{kg}$ (bien plus en pratique, pour chauffer les
cornues).
\textbf{21.} $\Delta_r H^\circ = \frac12(-226.8 + 946.2) = \qty{+360}{kJ}$ par
mole de zinc~; par tonne, $1.53 \times 10^4 \times 360 = \qty{5.5e6}{kJ}$, soit
\qty{5.5}{GJ} --- fournis en brûlant davantage de charbon autour des cornues.
\textbf{22.} Le carbone réduit l'oxyde de zinc dans les conditions standard au-dessus de
$\boldsymbol{T \approx \qty{1.22e3}{K}}$.
\end{solution}
